\runninghead{The Sealed and the Uncounted}
\properlane{sealed-number-uncounted-crowd}{One Church, Sealed and Uncounted}

The Lesson of this Mass counts the servants of God who are sealed, a hundred
and forty-four thousand, and then shows a multitude that no one could count.
This interpretation follows St Bede, who reads the number as the whole Church
and the seal as the sign of the cross, and St Augustine, who on the psalms of
this Mass bids the just rejoice in the Lord and not in themselves. On their
readings the saints whom the feast honours are those whom God has sealed and
gathered, from Israel and from every nation.

\subsection{The question: a number, and a multitude without number}
\label{sec:sealed-question}

John hears a number before he sees anything: \latin{centum quadraginta
quatuor millia signati, ex omni tribu filiorum Israel}, and the number is
then given again tribe by tribe, twelve thousand from each. Immediately after
it he sees \latin{turbam magnam, quam dinumerare nemo poterat}, from all
nations and tribes and peoples and tongues. Read as a count, the number
contradicts the multitude, or restricts the sealed to Israel and leaves the
nations outside the seal. The list of tribes adds a second difficulty. It is
not the list of Jacob's sons in the order of their birth: Juda stands first,
Dan is missing, and Manasses is counted beside his father Joseph. A third
difficulty is the seal itself. The angel forbids the four angels to hurt the
earth \latin{quoadusque signemus servos Dei nostri in frontibus eorum}, but
the Lesson does not say what the seal is, or what it does for those who
bear it.

\subsection{St Bede: the number is the whole Church}
\label{sec:sealed-bede-number}

St Bede, commenting on the Apocalypse verse by verse, takes the angel from
the rising of the sun to be the Lord incarnate, ``the day-spring from on
high'' (Lk 1:78), ``bearing the ensign of the cross, with which to seal His
own in their foreheads'' (\work{Explanatio Apocalypsis} I, on 7:2). The
command not to hurt the earth he reads of what followed the Passion, when the
dominion of the enemy and of worldly power was broken: ``to this end was the
empire of the nations broken up, that the face of the saints might be freely
marked with the seal of faith, which these had resisted'' (on 7:3). The seal
is the cross, and it is set where it can be seen.

On the number he is explicit: ``By this definite number is signified the
innumerable multitude of the whole Church, which is descended from the
patriarchs either by the lineage of nature, or the imitation of faith.'' His
warrant for counting the Gentiles among the tribes is St Paul's ``if ye be
Christ's, then are ye Abraham's seed'' (Gal 3:29). Twelve is multiplied by
twelve and raised to thousands, ``the cube of the number ten, by which is
represented the enduring life of the Church'' (on 7:4). The number is a
figure of completeness, and the completeness is that of the Church.

The order of the tribes he explains in the same way. John begins with Juda,
``from which tribe our Lord sprang,'' and leaves out Dan, from whom, ``as it is
said, Antichrist is to be born,'' citing Jacob's blessing, ``Let Dan become an
adder in the way'' (Gen 49:17). John ``has not proposed to set forth the order
of earthly generation, but according to the interpretation of the names the
powers of the Church, which from present confession and praise is hastening to
the right hand of eternal life.'' Juda means ``confession'' or ``praise,'' and
comes first, ``for before the first step of confession no one reaches the
height of good works''; Benjamin, ``the son of the right hand,'' comes last,
because after death is destroyed the elect are given the eternal inheritance
at Christ's right hand (on 7:5--8). Between them Bede reads each name as a
virtue in which the faithful make progress, from the works of Reuben to the
martyrdom of Zabulon and the unwearied prayer of Benjamin, and concludes that
``from the merits of individuals, should the most perfect beauty of the
Church, as the sum of a hundred and forty-four thousand, be made up.'' On his
reading the twelve thousand of each tribe are not twelve groups of people
but twelve kinds of holiness, each sealed ``in his own profession,'' which
together make up one Church.

\subsection{St Bede: the multitude before the throne}
\label{sec:sealed-bede-multitude}

At verse~9 Bede gives two readings of the multitude, and the study keeps
both. First, the vision ``announces the glory of those who are to overcome the
wickedness of the last persecution.'' Then, ``that which follows, `From all
nations, and tribes, and peoples, and tongues,' may also be thus understood,
that after enumerating the tribes of Israel, to whom the Gospel was first
preached, he desires to make mention of the salvation of the Gentiles as
well'' (on 7:9). Bede does not say that the number and the multitude are the
same people seen twice. He says that the number signifies the whole Church
(on 7:4), and, in his second reading, that the multitude adds the Gentiles to
the tribes. This interpretation joins the two and takes the number and the
multitude as one Church: the number as its completeness, the multitude as its
extent. That join is the study's, made from Bede's two comments.

The robes, he writes, signify baptism, and the palms ``the triumph of the
Cross,'' and ``in Christ they have overcome the world.'' The loud voice is
``great devotion, an unceasing praise'' (on 7:10). When the angels fall down
before the throne, ``In all the angels he has represented the persons of the
great multitude worshipping the Lord,'' though Bede allows that the verse may
also be read of the angelic spirits themselves (on 7:11). The sevenfold praise
of verse~12 is the Church's: ``The Church offers the sevenfold praise of
excellence unto the Lord, and in each of its members confesses to have
received this from Him'' (on 7:12).

Two of Bede's comments on verses that the Lesson does not read bear on whom
the feast honours. On the white robes washed in the blood of the Lamb (7:14)
he writes, ``He speaks not of martyrs alone. They are washed in their own
blood. But the blood of Jesus, the Son of God, cleanses the whole Church from
all sin.'' And after expounding the end of the chapter as the life to come, he
adds that the vision of the white-robed multitude ``may also be understood of
the present time, when `we are saved in hope''' (Rom 8:24; on 7:17). The feast
was first a feast of martyrs: Schuster traces it to a collective feast of
martyrs in fourth-century Syria and to the feast of the martyrs that Boniface
IV instituted when he consecrated the Pantheon as a church (\work{The
Sacramentary} V, pp.~208--209). On Bede's reading of
7:14 the white robes are not the martyrs' alone.

\subsection{St Augustine: rejoicing in the Lord}
\label{sec:sealed-augustine}

St Augustine does not comment on this chapter of the Apocalypse in the
works used here, and he says nothing of the seal or the number. He enters this
interpretation through the psalm verses that the Mass sings. On the Introit's
verse he writes, ``rejoice, O ye righteous, not in yourselves, for that is not
safe; but in the Lord. `For praise is comely to the upright': these praise the
Lord, who submit themselves unto the Lord'' (\work{En.\ in Ps.}\ 32 [NPNF
``Psalm XXXIII''], 1). The Introit antiphon bids all rejoice in the Lord on
the feast of all the saints, and the verse, read with Augustine, says in whom
the just rejoice: not in themselves.

Later in the same psalm, outside the verse the Mass sings, Augustine reads
``Blessed is the nation whose God is the Lord'' of ``one nation \dots\
belonging to the heavenly city,'' which ``not of itself, but by the gift of
God, hath been chosen'' (32:12; \S12). The Lesson's multitude is \latin{ex
omnibus gentibus}; to read Augustine's one nation beside it is this study's
join, not his. What his comment gives is the point this interpretation draws
from Bede as well: those who are counted are counted by God's choice and gift.

On the Gradual, Augustine is a second witness. Its verse sings that those who
seek the Lord shall lack no good, and he asks what good is meant. ``When thou
art filled with spiritual riches, canst thou be poor?'' The riches are ``such
jewelry of virtues, justice, truth, charity, faith, endurance,'' and of the
promise, ``what manner of good, I have already said'' (\work{En.\ in Ps.}\ 33
[NPNF ``Psalm XXXIV''], 14). The goods that the sealed are promised are, on
his reading, the virtues that Bede reads in the names of the tribes.

\subsection{The sealed in the other texts of the Mass}
\label{sec:sealed-elements}

The Introit and the Collect speak of the saints as one company. The antiphon
keeps the feast \latin{sub honore Sanctorum omnium}, and the Collect thanks
God for giving us to venerate the merits of all the saints \latin{sub una
\dots\ celebritate} and asks for mercy \latin{multiplicatis intercessoribus}.
Durandus explains the one celebration by the saints' number: so many could
not each be given a feast (\work{Rationale} VII.34). The Gradual calls them
\latin{Sancti eius}, his saints, and names them by their fear of the Lord; in
the psalm's preceding verse, which the Mass does not sing, ``the angel of the
Lord shall encamp round about them that fear him: and shall deliver them''
(Ps 33:8). To read that angel with the Lesson's angel who holds back the four
angels from the earth is the study's comparison, not the psalm's or Bede's.

The Offertory sings of the just in death, and the Postcommunion of the
faithful on earth. The souls of the just are in the hand of God and ``seemed
to die''; Bede's seal is set on those who ``in the midst of adverse things are
faithful confessors of His Name'' (on 7:14). The Postcommunion asks that the
\latin{fideles populi} be defended by the saints' unceasing prayer. Schuster
reads those peoples as the Church on earth, ``journeying towards eternity
with the eyes, and in the light of faith'' (p.~211); on this interpretation
they are the sealed who have not yet come to stand before the throne. The
Sunday's Collect, said after the feast's, asks God to loose the sins of his
\latin{populi}, and its Postcommunion asks that those who share in divine
things be not left subject to human dangers. These prayers belong to the
Sunday and say nothing of the seal; the study notes only that the peoples who
ask for the saints' prayer also ask to be absolved.

\subsection{The strongest difficulty: a number that limits}
\label{sec:sealed-difficulty}

The strongest objection is the plain sense of the Lesson's own words. The
number is exact, the tribes are Israel's, and John says he \latin{heard} the
number of the sealed before he \latin{saw} the multitude; a reading that takes
the number for the whole Church may seem to drop \latin{centum quadraginta
quatuor millia} in order to reach the crowd. Bede's answer is to keep the
number and to read it, giving his reasons: Paul's seed of Abraham by faith,
the arithmetic of twelve and ten, and the order of the tribes that is not the
order of birth. He does not explain the number away. The tribes and the
twelve thousands remain in his reading, as the powers of the Church from
confession to the right hand.

A different reading of the same verses is held by Lucien Fromage, who
continued Dom Prosper Guéranger's \work{Liturgical Year}. He takes the Lesson
as a second census. At his first coming the God-Man, using Caesar Augustus,
had the world enrolled, and now, Fromage writes, \latin{l'heure a sonné d'un
autre recensement, qui doit consigner au livre de vie le résultat des
opérations du salut}: the hour has struck for another census, which is to
record in the book of life the outcome of the work of salvation. He cites St
Gregory the Great's Christmas homily on the census, and adds that because
many withdrew by their own fault from the first enrolment, which comprised
all men in the redemption, a second and definitive one was needed to cut the
guilty off from it (\work{L'Année liturgique}, vol.~15, p.~94; Gregory,
\work{Hom.\ in Ev.}\ 8). Bede reads the number as the whole Church made up of
its members' virtues; Fromage reads the sealing as the last reckoning, which
enrols the saved and strikes out the rest. The two readings are not the same,
and the study does not merge them. On Fromage's reading the multitude is
still uncounted, but the census that precedes it is also a judgment.

Schuster draws the opposite conclusion from the size of the multitude: ``Not
only the hundred and forty-four thousand predestined sons of Abraham, but a
turba magna of every age, sex, and condition'' enter heaven, and ``It is not
then so hard to be saved, as St John writes that he could not count the number
of the elect'' (p.~210). He adds at once one condition, the seal of the Tau
that Ezechiel saw set on the foreheads of those who sigh and mourn, and he
reads it with St Paul's ``as you are partakers of the sufferings, so shall you
be also of the consolation'' (2~Cor 1:7). Schuster's 144,000 are sons of
Abraham, and his multitude adds others to them; Bede's 144,000 are already
the whole Church. On none of the three readings is the seal given to those
who do not bear it.

\begin{fourSenses}
\item[Literal.] John sees an angel from the east forbid the angels of
destruction to hurt the earth until the servants of God are sealed on their
foreheads; he hears their number, twelve thousand from each of twelve tribes;
he then sees an uncounted multitude from every nation in white robes with
palms before the throne and the Lamb, and the angels' sevenfold praise.
\item[Allegorical.] The angel is Christ incarnate, the seal his cross, and
the hundred and forty-four thousand the whole Church, descended from the
patriarchs by nature or by faith; the multitude is the victors over the last
persecution or, in Bede's second reading, the Gentiles gathered after Israel,
washed in the blood of the Lamb (Bede).
\item[Moral.] The seal is the faith confessed openly on the forehead, and the
tribes are the virtues of the faithful, from confession to perseverance in
prayer (Bede). The just rejoice in the Lord and not in themselves, and the
goods they are promised are justice, truth, charity, faith and endurance
(Augustine).
\item[Anagogical.] The elect stand at last at the right hand of Christ, which
Benjamin's place at the end of the list signifies, and give God unceasing
praise; the vision is also true now of those who are saved in hope (Bede).
\end{fourSenses}
